\chapter{Razprava in sklepne ugotovitve}
\label{ch:razprava-sklep}

\section{Odgovori na raziskovalna vprašanja}

Na prvo raziskovalno vprašanje lahko odgovorimo pritrdilno. Kvantizirani model
YAMNet-1024 se na mikrokrmilniku STM32N6 izvaja sproti. Predobdelava enega
bloka traja povprečno \SI{1.43}{\milli\second}, izvajanje nevronske mreže
\SI{9.38}{\milli\second}, skupaj pa \SI{10.81}{\milli\second}. Nov blok se
pripravi vsakih \SI{480}{\milli\second}, zato računanje porabi približno
\SI{2.25}{\percent} tega intervala. Tudi aplikacija, aktivacije, uteži in dva
zaslonska medpomnilnika ostanejo znotraj svojih pomnilniških območij. To
potrjuje tehnično izvedljivost celotne lokalne poti, ne pomeni pa, da je čas od
nastanka zvoka do opozorila samo \SI{10.81}{\milli\second}, saj mora sistem
najprej zbrati vhodni blok in po potrebi potrditi več zaporednih odločitev.

Na drugo vprašanje odgovarja neodvisni končni fizični preizkus. Sistem je
pričakovani razred vsaj enkrat potrdil v 57 od 60 poskusov oziroma v
\SI{95.0}{\percent} poskusov. Pri štirih nevarnostnih razredih je bil uspešen
v 39 od 40 poskusov. Strožje merilo prevladujočega izhoda je bilo pravilno v
54 od 60 poskusov. Govor ni sprožil nobenega nevarnostnega alarma, pri razredu
ostalo pa se je nevarnostni izhod pojavil v šestih od desetih poskusov. Sistem
je torej v preskusni postavitvi dobro zaznaval ciljne dogodke, zavračanje
nekaterih kratkih vsakdanjih zvokov pa ostaja njegova glavna slabost.

Tretje vprašanje nima enega samega odgovora za vse razrede. Pasji lajež,
razbitje stekla in strel so bili v končnem preizkusu potrjeni v vseh desetih
poskusih. Sirena, govor in ostalo so dosegli devet pravilnih potrditev od
desetih. Rezultati neuspešnih poskusov kažejo različne vzroke: ocena govora in
sirene je ostala pod razrednim pragom, kratek pok pa je bil podoben razbitju
stekla. Na delovanje zato vplivajo kakovost in glasnost vira, položaj dogodka v
skoraj enosekundnem vhodnem bloku, podobnost med razredi ter izbrani pragovi in časovna
pravila.

\section{Razprava o uporabnosti rezultatov}

Prototip je namenjen pretvarjanju izbranih nevarnih zvokov v vidno opozorilo
za gluhe in naglušne uporabnike. Njegova praktična prednost je samostojno
delovanje: mikrofon, predobdelava, model, odločanje in prikaz so na isti
napravi. Internetna povezava, oddaljeni strežnik in uporabniški račun niso
potrebni. Običajna programska oprema surovega zvoka ne shranjuje, začasni
vzorci pa se po sprotni obdelavi prepišejo. To zmanjša izpostavljenost zvoka iz
domačega okolja in hkrati omogoča delovanje tudi ob izpadu omrežja.

Zaslon jasno razlikuje med čakanjem, nezanesljivo napovedjo, običajnim zvokom
in potrjenim nevarnostnim dogodkom. Takšna izvedba je primerna kot raziskovalni
prototip in kot osnova za nadaljnji razvoj namenskega opozorilnega pripomočka.
Ni pa nadomestilo za certificiran požarni, protivlomni ali medicinski sistem.
Preizkus je obsegal nadzorovano predvajanje v enem prostoru, zato iz rezultata
ne moremo sklepati, da bo naprava enako zanesljiva pri vseh resničnih zvokih,
razdaljah in okoljih.

\section{Omejitve projekta}

V vsakem sistemskem razredu je bilo deset preskusnih posnetkov, zato so
razredna območja zaupanja široka. Zvoki so bili predvajani prek enega zvočnika
na razdalji \SI{30}{\centi\meter} v enem prostoru. Tak postopek zagotavlja
ponovljivost, ne zajame pa odmeva, oddaljenosti, prekrivanja več zvokov in
različnih mikrofonov pri vsakdanji uporabi. Uporabljeni posnetki so bili novi
glede na učenje in izbiro pragov, vendar še vedno izvirajo iz javnih zbirk in
ne iz daljšega terenskega snemanja ciljnega okolja.

Sistem pozna omejen nabor nevarnosti. Posebej zahtevni so kratki dogodki, ki
so si spektralno podobni. Razred ostalo je raznolik in ga 630 učnih virov ne
more izčrpno opisati; zato so pok, aplavz in nekateri gospodinjski zvoki še
vedno povzročili lažne nevarnostne izhode. Izmerjeni čas temelji na notranji
telemetriji, ne na zunanjem električnem signalu, poraba energije pa ni bila
izmerjena. Prav tako uporabniška študija z gluhimi ali naglušnimi osebami ni
bila izvedena. Namen za ciljno skupino je zato zasnovna motivacija, ne dokaz
klinične ali vsakdanje uporabnosti.

\section{Možnosti nadaljnjega dela}

Najpomembnejši naslednji korak je razširitev razreda ostalo z lokalnimi
gospodinjskimi zvoki ter s težkimi negativnimi primeri, pri katerih trenutni
sistem napove nevarnost. Temu naj sledi novo učenje in popolnoma ločen končni
preskus. Smiselni so preskusi na več razdaljah, v več prostorih, ob hrupu in ob
hkratnem pojavu govora ter nevarnostnega dogodka. Lasten uravnotežen nabor bi
bolje opisal pričakovano okolje kot samo javni posnetki.

Celotno zakasnitev bi bilo mogoče natančneje izmeriti z zunanjim signalom in
osciloskopom, porabo pa z merilnikom energije. Za uporabniški del bi bilo treba
skupaj z gluhimi in naglušnimi osebami preveriti razumljivost opozoril,
velikost besedila, barve ter morebitno dodajanje vibracijskega ali svetlobnega
izhoda. Morebitno razširjanje nabora nevarnosti naj temelji na dovolj kakovostnih
učnih podatkih in ločenem vrednotenju, ne le na dodajanju novih oznak.

\section{Sklepne ugotovitve}

V diplomskem delu je bil razvit samostojen sistem za zaznavanje nevarnih
zvočnih dogodkov na razvojni plošči STM32N6570-DK. Vgrajeni mikrofon
neprekinjeno zajema zvok, lastna predobdelava pripravi strnjeni spektrogram,
kvantizirani model YAMNet-1024 se izvaja z nevronskim
pospeševalnikom Neural-ART, odločitveni filter pa rezultat pretvori v stabilno
vidno opozorilo. Celotna obdelava poteka lokalno, brez dostopa do računalniškega
oblaka in brez trajnega shranjevanja surovega zvoka v običajnem delovanju.

Razvoj je pokazal, da pri vgrajenem razpoznavanju ni dovolj samo namestiti
prednaučenega modela. Za kratke dogodke so bili potrebni namensko izrezovanje,
povečanje podatkov, razred ostalo, razredni pragovi in časovno filtriranje.
Končna različica je v neodvisnem fizičnem preizkusu pravilno potrdila 57 od 60
poskusov, pri nevarnostnih razredih pa 39 od 40. Obdelava enega bloka je
trajala \SI{10.81}{\milli\second} in je ostala znotraj razpoložljivega
pomnilnika. Preostali lažni alarmi pri razredu ostalo jasno določajo naslednjo
razvojno nalogo. Glavni rezultat dela je zato delujoč in merljivo dovolj hiter
lokalni prototip ter sledljiv postopek, ki njegove prednosti loči od njegovih
še odprtih omejitev.
